%------------------------------------------------------------------------------%

\usepackage{calculo_varias_variaveis-1}

\title {Máximos e Mínimos em Regiões Fechadas}

%------------------------------------------------------------------------------%
\begin{document}

\addtitle

%------------------------------------------------------------------------------%
\section{Extremos Absolutos em Regiões Fechadas e Limitadas}

%------------------------------------------------------------------------------%
%\framedgraphic{Derivada Direcional}{derivada_direcional}{}

%------------------------------------------------------------------------------%
\begin{frame}{Região Fechada e Limitada}
  \begin{description}[<+->]
    \item[Fechada] Contém todos os pontos de fronteira
    \item[Limitada] Está dentro de um círculo centrado na origem
  \end{description}
\end{frame}

%------------------------------------------------------------------------------%
\begin{frame}{Função Contínua}
  Uma função contínua em uma região fechada e limitada possui um máximo e um
  mínimo absolutos (globais)
\end{frame}

%------------------------------------------------------------------------------%
\begin{frame}{Encontando os Extremos Absolutos}
  Encontando os Extremos Absolutos de uma função $f$ \\
  em uma Região Fechada e Limitada $R$
  \pause
  \begin{enumerate}[<+->]
    \item Listar os pontos críticos de $f$ no interior de $R$
    \item Listar os pontos onde $f$ pode ter máximos ou mínimos na fronteira de $R$
    \item Calcular o valor de $f$ nos pontos listados
    \item Encontrar os valores máximo e mínino nas listas
  \end{enumerate}
\end{frame}

%------------------------------------------------------------------------------%
\addtocounter{exemplo}{1}
\begin{frame}{Exemplo \theexemplo}
  Encontre os valores máximo e mínimo absolutos de
  \[
    f(x, y) = 2 + 2x + 2y - x^2 - y^2
  \]
  na região triangular no primeiro quadrante limitada pelas retas \\
  \(x=0\), \(y=0\) e \(y = 9-x\)
\end{frame}

\begin{frame}{Exemplo \theexemplo}
  Pontos críticos interiores
  \[
    f_x(x, y)
    \pause
    = \D{}{x}\left(2 + 2x + 2y - x^2 - y^2\right)
    \pause
    = 2-2x
    \pause
    = 0
    \pause
    \quad\Rightarrow\quad
    x = 1
  \]
  \[
    \pause
    f_y(x, y)
    \pause
    = \D{}{y}\left(2 + 2x + 2y - x^2 - y^2\right)
    \pause
    = 2-2y
    \pause
    = 0
    \pause
    \quad\Rightarrow\quad
    y = 1
  \]
  \pause
  É um ponto interior

  \pause
  \vspace{\baselineskip}
  Valor da função
  \[
    f(1, 1)
    \pause
    = \left(2 + 2x + 2y - x^2 - y^2\right)\evalat{(1, 1)}{}
    \pause
    = 2 + 2 + 2 - 1 - 1
    \pause
    = 4
  \]
\end{frame}

\begin{frame}{Exemplo \theexemplo}
  Pontos na fronteira

  \pause
  \vspace{\baselineskip}
  Vamos tratar cada segmento em separado
  \begin{enumerate}[<+->]
    \item \(y=0\), \(0\leq x\leq 9\)
    \item \(x=0\), \(0\leq y\leq 9\)
    \item \(y=9-x\), \(0\leq x\leq 9\)
  \end{enumerate}
\end{frame}

\begin{frame}{Exemplo \theexemplo}
  Segmento \(y=0\), \(0\leq x\leq 9\)

  \pause
  \vspace{\baselineskip}
  Função de uma variável, $x$, no intervalo fechado \(0\leq x\leq 9\)
  \[
    \pause
    g(x)
    \pause
    = f(x, 0)
    \pause
    = 2 + 2x - x^2
  \]

  \pause
  \vspace{0.3\baselineskip}
  Pontos interiores
  \[
    g'(x) = 2 - 2x = 0
    \pause
    \quad\Rightarrow\quad
    x = 1
  \]
  \pause
  Extremos \(x=0\) e \(x=9\)

  \pause
  \vspace{0.5\baselineskip}
  Pontos no plano
  \((0, 0)\), \((1, 0)\), \((9, 0)\)
\end{frame}

\begin{frame}{Exemplo \theexemplo}
  Valores da função
  \[
    \pause
    f(0, 0)
    \pause
    = g(0)
    \pause
    = 2
  \]
  \[
    \pause
    f(1, 0)
    \pause
    = g(1)
    \pause
    = 3
  \]
  \[
    \pause
    f(9, 0)
    \pause
    = g(9)
    \pause
    = 2 + 2\times 9 - 9^2
    \pause
    = -61
  \]
\end{frame}

\begin{frame}{Exemplo \theexemplo}
  Segmento \(x=0\), \(0\leq y\leq 9\)

  \pause
  \vspace{\baselineskip}
  Função de uma variável, $y$, no intervalo fechado \(0\leq y\leq 9\)
  \[
    \pause
    h(y)
    \pause
    = f(0, y)
    \pause
    = 2 + 2y - y^2
  \]

  \pause
  \vspace{0.3\baselineskip}
  Pontos interiores
  \[
    \pause
    h'(y) = 2 - 2y = 0
    \pause
    \quad\Rightarrow\quad
    y = 1
  \]
  \pause
  Extremos \(y=0\) e \(y=9\)

  \pause
  \vspace{0.5\baselineskip}
  Pontos no plano
  \((0, 0)\), \((0, 1)\), \((0, 9)\)
\end{frame}

\begin{frame}{Exemplo \theexemplo}
  Valores da função
  \[
    \pause
    f(0, 0)
    \pause
    = h(0)
    \pause
    = 2
  \]
  \[
    \pause
    f(0, 1)
    \pause
    = h(1)
    \pause
    = 3
  \]
  \[
    \pause
    f(0, 9)
    \pause
    = g(9) \pause
    = 2 + 2\times 9 - 9^2
    \pause
    = -61
  \]
\end{frame}

\begin{frame}{Exemplo \theexemplo}
  Segmento \(y=9-x\), \(0\leq x\leq 9\)

  \pause
  \vspace{\baselineskip}
  Função de uma variável, $x$, no intervalo fechado \(0\leq y\leq 9\)
  \[
    \pause
    p(x)
    \pause
    = f(x, 9-x)
    \pause
    = 2 + 2x + 2(9-x) - x^2 - (9-x)^2
%    \pause
%    = 2 + 2x + 2\times 9 -2x - x^2 - 9^2 +2\times 9x - x^2
  \]
  \[
    \hphantom{p(x)}
    \pause
    = 2 + 18 - 81 + 2x -2x +18x -x^2 - x^2
  \]
  \[
    \hphantom{p(x)}
    \pause
    = -61 + 18x -2x^2
  \]

  \pause
  \vspace{0.3\baselineskip}
  Pontos interiores
  \[
    \pause
    p'(x) = 18 - 4x = 0
    \pause
    \quad\Rightarrow\quad
    x = \frac{18}{4}
    \pause
    = \frac{9}{2}
  \]
  \pause
  Extremos \(x=0\) e \(x=9\)

  \pause
  \vspace{0.5\baselineskip}
  Pontos no plano
  \((0, 9)\), \(\left(\sfrac{9}{2}, \sfrac{9}{2}\right)\), \((9, 0)\)
\end{frame}

\begin{frame}{Exemplo \theexemplo}
    Valores da função
    \[
      \pause
      f(0, 9) = -61
    \]
    \[
      \pause
      f\left(\frac{9}{2}, \frac{9}{2}\right)
      \pause
      = 2 + 2\frac{9}{2} + 2\frac{9}{2} - \left(\frac{9}{2}\right)^2 - \left(\frac{9}{2}\right)^2
      \pause
      = 2 + 9 + 9 -\frac{81}{4} -\frac{81}{4}
      \]
      \[
      \pause
      \hphantom{p(x)}
      = 20 -\frac{81}{2}
      \pause
      = \frac{40-81}{2}
      \pause
      = \frac{-41}{2}
    \]
    \[
      \pause
      f(9, 0) = -61
    \]
\end{frame}

\begin{frame}{Exemplo \theexemplo}
  Máximo absoluto de $f$ em $R$ é 4 e ocorre no ponto \((1, 1)\)

  \vspace{\baselineskip}
  Mínimo absoluto de $f$ em $R$ é $-61$ e ocorre nos pontos \((9, 0)\) e \((0, 9)\)
\end{frame}

%------------------------------------------------------------------------------%
\addtocounter{exemplo}{1}
\begin{frame}{Exemplo \theexemplo}
  Encontre os máximos e mínimos absolutos da função
  \[
    f(x, y) = x^2 + 2y^2 - x
  \]
  na região
  \[
    x^2 + y^2 \leq 1
  \]
\end{frame}

\begin{frame}{Exemplo \theexemplo}
  Pontos críticos interiores
  \[
    f_x(x, y)
    \pause
    = \D{}{x}\left(x^2 + 2y^2 - x\right)
    \pause
    = 2x-1
    \pause
    = 0
    \pause
    \quad\Rightarrow\quad
    x = \frac{1}{2}
  \]
  \[
    \pause
    f_y(x, y)
    \pause
    = \D{}{x}\left(x^2 + 2y^2 - x\right)
    \pause
    = 4y
    \pause
    = 0
    \pause
    \quad\Rightarrow\quad
    y = 0
  \]
  \pause
  É um ponto interior

  \pause
  \vspace{\baselineskip}
  Valor da função
  \[
    f\left(\sfrac{1}{2}, 0\right)
    \pause
    = \left(x^2 + 2y^2 - x\right)\evalat{\left(\sfrac{1}{2}, 0\right)}{}
    \pause
    = \left(\frac{1}{2}\right)^2 + 2\times 0 ^2 - \frac{1}{2}
    \pause
    = \frac{1}{4} - \frac{2}{4}
    \pause
    = -\frac{1}{4}
  \]
\end{frame}

\begin{frame}{Exemplo \theexemplo}
  Fronteira
  \[
    x^2+y^2 = 1
    \pause
    \quad\Rightarrow\quad
    y^2 = 1 - x^2
  \]
  \pause
  com \(-1\leq x \leq 1\)

  \pause
  \vspace{\baselineskip}
  Função $f$ restrita à fronteira
  \[
    g(x)
    \pause
    = x^2 + 2\left(1 - x^2\right) -x
    \pause
    = x^2 + 2 - 2x^2 - x
    \pause
    = 2 - x^2 - x
  \]
\end{frame}

\begin{frame}{Exemplo \theexemplo}
  Candidatos a máximos e mínimos na fronteira

  \pause
  \vspace{\baselineskip}
  Extremidade esquerda
  \[
    x = -1
    \pause
    \quad\Rightarrow\quad
    \pause
    y^2 = 1 - (-1)^2 = 0
  \]
  \pause
  Extremidade direita
  \[
    \pause
    x = 1
    \pause
    \quad\Rightarrow\quad
    y^2 = 1 - 1^2
    \pause
    = 0
  \]
\end{frame}

\begin{frame}{Exemplo \theexemplo}
  Interior
  \[
    \pause
    g'(x)
    \pause
    = \frac{d}{dx}\left(2 - x^2 - x\right)
    \pause
    = -2x - 1
    \pause
    = 0
    \pause
    \quad\Rightarrow\quad
    x = -\frac{1}{2}
  \]
  \[
    \pause
    y^2 = 1 - x^2
    \pause
    = 1 - \left(-\frac{1}{2}\right)^2
    \pause
    = \frac{4}{4} - \frac{1}{4}
    \pause
    = \frac{3}{4}
 \]
 \[
    \pause
    \abs{y} = \sqrt{\frac{3}{4}}
  \]
  \[
    \pause
    y
    = \pm\sqrt{\frac{3}{4}}
    \pause
    = \pm\frac{\sqrt{3}}{2}
  \]
\end{frame}

\begin{frame}{Exemplo \theexemplo}
  Avaliando a função nos pontos
  \[
    \pause
    f(-1, 0)
    \pause
    = \left(x^2 + 2y^2 - x\right)\evalat{(-1, 0)}{}
    \pause
    = 1 + 0 - (-1)
    \pause
    = 2
  \]
  \[
    \pause
    f(1, 0)
    \pause
    = \left(x^2 + 2y^2 - x\right)\evalat{(1, 0)}{}
    \pause
    = 1 + 0 - 1
    \pause
    = 0
  \]
\end{frame}

\begin{frame}{Exemplo \theexemplo}
  \[
    f\left(-\frac{1}{2}, \frac{\sqrt{3}}{2}\right)
    \pause
    = \left(x^2 + 2y^2 - x\right)\evalat{\left(-\frac{1}{2}, \frac{\sqrt{3}}{2}\right)}{}
    \pause
    = \left(-\frac{1}{2}\right)^2 + 2\left(\frac{\sqrt{3}}{2}\right)^2 - \left(-\frac{1}{2}\right)
  \]
  \[
    \pause
    \hphantom{f\left(-\frac{1}{2}, \frac{\sqrt{3}}{2}\right)}\,
    = \frac{1}{4} + 2\frac{3}{4} + \frac{2}{4}
    \pause
    = \frac{9}{4}
    = 2,25
  \]

  \pause
  \[
    f\left(-\frac{1}{2}, -\frac{\sqrt{3}}{2}\right)
    \pause
    = \left(x^2 + 2y^2 - x\right)\evalat{\left(-\frac{1}{2}, -\frac{\sqrt{3}}{2}\right)}{}
    \pause
    = \left(-\frac{1}{2}\right)^2 + 2\left(-\frac{\sqrt{3}}{2}\right)^2 - \left(-\frac{1}{2}\right)
  \]
  \[
    \pause
    \hphantom{f\left(-\frac{1}{2}, -\frac{\sqrt{3}}{2}\right)}\,
    = \frac{1}{4} + 2\frac{3}{4} + \frac{2}{4}
    \pause
    = \frac{9}{4}
    = 2,25
  \]
\end{frame}

\begin{frame}{Exemplo \theexemplo}
  Máximo absoluto de $f$ na região é \(\frac{9}{4}\) e ocorre nos pontos \\
  \(\left(-\frac{1}{2}, -\frac{\sqrt{3}}{2}\right)\) e
  \(\left(-\frac{1}{2},  \frac{\sqrt{3}}{2}\right)\)

  \vspace{\baselineskip}
  Mínimo absoluto de $f$ na região é \(-\frac{1}{4}\)
  e ocorre no ponto
  \(\left(\frac{1}{2}, 0\right)\)
\end{frame}

%------------------------------------------------------------------------------%
\section{Lista Mínima}

%------------------------------------------------------------------------------%
\begin{frame}{Lista Mínima}
  Cálculo Vol. 2 do Thomas 12\textsuperscript{a} ed. -- Seção 14.7
  \begin{enumerate}
    \item Estudar o texto da seção
    \item Resolver os exercícios: 31, 34, 37, 38
  \end{enumerate}

  \vfill
  Atenção: A prova é baseada no livro, não nas apresentações
\end{frame}

%------------------------------------------------------------------------------%
\end{document}
%------------------------------------------------------------------------------%
